\chapter*{Introduction générale}
\addcontentsline{toc}{chapter}{Introduction générale}

La transformation numérique des organisations ne se limite plus à la dématérialisation documentaire. Elle implique désormais une capacité à modéliser les processus métier, à centraliser l'information opérationnelle, à tracer les interactions, à sécuriser les accès et à produire des indicateurs permettant une prise de décision rapide et fiable. Dans ce contexte, la gestion du parc informatique représente un enjeu stratégique pour toute institution, car elle conditionne à la fois la disponibilité des outils de travail, la qualité du support, la maîtrise des coûts, la conformité des opérations et la satisfaction des utilisateurs finaux.

Dans de nombreuses structures académiques et professionnelles, la gestion des équipements reste encore fragmentée entre plusieurs fichiers, échanges informels, validations manuelles et outils partiellement interconnectés. Cette dispersion entraîne plusieurs difficultés récurrentes : absence de visibilité globale sur l'état réel des actifs, complexité de suivi des affectations, délais dans le traitement des incidents, manque de traçabilité lors des retours et insuffisance des tableaux de bord pour piloter la maintenance ou anticiper les besoins. La qualité de service dépend alors fortement d'interventions humaines répétitives, peu industrialisées et difficilement auditables.

Le projet présenté dans ce mémoire répond précisément à cette problématique par la mise en œuvre d'une plateforme web intelligente de gestion du parc informatique et des affectations matérielles. L'ambition n'est pas seulement de construire une application fonctionnelle, mais de proposer une solution cohérente, évolutive et académiquement défendable, s'appuyant sur des choix techniques modernes et sur une modélisation rigoureuse des besoins. La plateforme vise à couvrir l'authentification, la gestion des rôles, l'inventaire des équipements, le cycle de vie des affectations, le suivi des maintenances, les notifications en temps réel, l'audit des opérations et la génération de rapports.

Sur le plan technologique, le projet adopte une architecture full-stack articulée autour d'un frontend React pour l'expérience utilisateur, d'un backend Spring Boot pour l'exposition des services métiers, d'une base PostgreSQL pour la persistance relationnelle, de mécanismes de sécurité basés sur JWT, d'un canal WebSocket pour les interactions temps réel et d'un déploiement conteneurisé à l'aide de Docker. Ce choix d'architecture répond aux exigences contemporaines de modularité, de performance, de maintenabilité et d'intégration continue, tout en restant accessible dans un cadre pédagogique de projet de fin d'année.

L'intérêt académique du sujet réside dans le croisement de plusieurs dimensions complémentaires. Premièrement, il mobilise des compétences d'analyse fonctionnelle et de modélisation UML, indispensables à toute démarche d'ingénierie logicielle structurée \cite{booch2005}. Deuxièmement, il impose des arbitrages d'architecture logicielle et de sécurité applicative, notamment autour des API REST, du contrôle d'accès, de la séparation des responsabilités et de la traçabilité \cite{fielding2000,owasp2021}. Troisièmement, il met en évidence l'importance de l'expérience utilisateur et de la lisibilité visuelle dans les applications métier, souvent négligées alors qu'elles constituent un facteur majeur d'adoption.

Le présent rapport a donc pour objectif de documenter de manière détaillée l'ensemble du processus ayant conduit à la réalisation de la plateforme. Il expose le contexte général du projet, l'analyse des besoins, la modélisation, l'architecture technique, la conception de la base de données, la réalisation logicielle, les considérations UI/UX, la stratégie de tests, le déploiement ainsi que les résultats obtenus. Le mémoire se veut à la fois descriptif, analytique et critique, dans l'esprit d'un véritable rapport d'ingénierie.

La suite du document est organisée comme suit :
\begin{itemize}
  \item le premier chapitre présente le contexte général, la problématique et les objectifs du projet ;
  \item le deuxième chapitre détaille l'analyse fonctionnelle et les exigences ;
  \item le troisième chapitre expose la modélisation UML de la solution ;
  \item les chapitres quatre, cinq et six décrivent respectivement l'architecture, la base de données et la réalisation technique ;
  \item le septième chapitre est consacré au design UI/UX et aux captures d'écran ;
  \item les chapitres huit et neuf portent sur les tests et le déploiement ;
  \item le dixième chapitre analyse les résultats, les limites et les perspectives ;
  \item enfin, la conclusion générale synthétise les apports du projet avant les annexes techniques.
\end{itemize}

\chapter{Contexte général du projet}
\label{chap:contexte}

\section{Présentation du contexte}
Le parc informatique d'une organisation regroupe l'ensemble des ressources matérielles et logicielles nécessaires à son fonctionnement quotidien : ordinateurs, périphériques, terminaux mobiles, accessoires, espaces de stockage, équipements de remplacement, historiques de maintenance et mouvements d'affectation. Dans un établissement d'enseignement supérieur ou dans une structure à forte composante numérique, ces actifs ne constituent pas de simples objets inventoriables ; ils représentent un capital opérationnel critique dont la disponibilité influence directement la productivité des équipes, la qualité des services et la continuité des activités pédagogiques ou administratives.

Dans ce cadre, la gestion du parc informatique doit répondre à plusieurs impératifs simultanés :
\begin{itemize}
  \item disposer d'un inventaire centralisé, à jour et interrogeable ;
  \item suivre précisément qui utilise quel matériel, dans quel service et pour quelle durée ;
  \item détecter les équipements inactifs, obsolètes ou en risque de panne ;
  \item planifier et tracer les interventions de maintenance ;
  \item sécuriser l'accès aux données et conserver l'historique des opérations.
\end{itemize}

Le projet s'inscrit donc dans une logique de professionnalisation des processus IT. Il répond à une attente forte de rationalisation et de pilotage, en proposant une plateforme capable de transformer des opérations dispersées en un système cohérent, sécurisé et orienté décision.

\section{Présentation du besoin}
Le besoin initial observé peut être résumé par trois constats majeurs. D'abord, la dispersion de l'information entre différents supports rend difficile l'obtention d'une vision globale du parc. Ensuite, les processus d'affectation et de retour sont souvent peu formalisés, ce qui génère des ambiguïtés sur la responsabilité d'un équipement, son état réel et son emplacement. Enfin, les activités de maintenance et de support manquent parfois d'outils de suivi consolidés, notamment lorsqu'il s'agit de prioriser les interventions ou d'analyser les tendances.

La plateforme développée vise à couvrir ce besoin de façon transversale. Elle fournit un espace unique dans lequel les utilisateurs habilités peuvent consulter un tableau de bord synthétique, administrer les équipements, gérer les utilisateurs, créer des affectations, visualiser les incidents et les maintenances, recevoir des notifications et produire des rapports. Elle s'appuie sur un modèle de données métier structuré et sur des composants de sécurité adaptés à une application d'entreprise.

\section{Problématique}
La problématique centrale du projet peut être formulée comme suit :

\begin{keybox}[Problématique]
Comment concevoir et implémenter une plateforme web moderne, sécurisée et évolutive permettant de centraliser la gestion du parc informatique, de fiabiliser les affectations matérielles, de tracer les opérations critiques et d'améliorer le pilotage opérationnel grâce à une interface claire et à des services métier cohérents ?
\end{keybox}

Cette problématique se décline en sous-questions :
\begin{itemize}
  \item comment modéliser le cycle de vie des équipements et des affectations sans alourdir l'expérience utilisateur ;
  \item comment garantir la sécurité des accès et la traçabilité des actions ;
  \item comment articuler de manière fluide frontend, backend, base relationnelle et notifications temps réel ;
  \item comment produire un système à la fois convaincant sur le plan pédagogique et crédible sur le plan professionnel.
\end{itemize}

\section{Objectifs du projet}
\subsection{Objectif général}
L'objectif général consiste à développer une plateforme web intelligente destinée à la gestion centralisée du parc informatique et des affectations matérielles, avec un haut niveau de lisibilité fonctionnelle, une sécurité adaptée à un environnement multi-rôles et une architecture technique capable d'évoluer vers un contexte de production.

\subsection{Objectifs fonctionnels}
Les objectifs fonctionnels retenus pour le projet sont synthétisés dans le tableau \ref{tab:objectifs-fonctionnels}.

\begin{table}[H]
\centering
\caption{Objectifs fonctionnels du projet}
\label{tab:objectifs-fonctionnels}
\begin{tabularx}{\textwidth}{>{\bfseries}p{3.4cm}X}
\toprule
Domaine & Objectif visé \\
\midrule
Authentification & Permettre une connexion sécurisée, la gestion des rôles et la persistance de session. \\
Inventaire & Centraliser la fiche de chaque équipement, son statut, son code d'inventaire et son historique. \\
Affectations & Formaliser les demandes, validations, affectations actives et restitutions. \\
Maintenance & Suivre les interventions, les incidents, les priorités et les états d'avancement. \\
Notifications & Informer les utilisateurs des événements clés en temps réel. \\
Reporting & Produire des exports et synthèses exploitables pour le pilotage. \\
Audit & Conserver la trace des actions sensibles sur les entités critiques. \\
\bottomrule
\end{tabularx}
\end{table}

\subsection{Objectifs non fonctionnels}
Au-delà des fonctions visibles, la plateforme doit respecter plusieurs objectifs non fonctionnels essentiels :
\begin{itemize}
  \item offrir une interface ergonomique, responsive et professionnelle ;
  \item garantir la confidentialité des accès via un mécanisme d'authentification robuste ;
  \item assurer la maintenabilité grâce à une architecture modulaire ;
  \item permettre l'extension progressive du périmètre fonctionnel ;
  \item rester déployable dans différents environnements au moyen de Docker \cite{merkel2014docker,dockerdocs}.
\end{itemize}

\section{Périmètre du projet}
Le périmètre fonctionnel de la version développée couvre les modules suivants :
\begin{enumerate}
  \item gestion des comptes utilisateurs, des rôles et des permissions ;
  \item gestion des équipements informatiques et de leurs attributs principaux ;
  \item gestion des affectations et des restitutions ;
  \item tableau de bord analytique et flux d'activité récents ;
  \item maintenance, tickets, historique et notifications ;
  \item génération de rapports PDF et Excel ;
  \item intégration de services transverses : audit, sécurité JWT, WebSocket, déploiement Docker.
\end{enumerate}

Le projet ne couvre pas, dans cette itération, certaines fonctionnalités industrielles avancées telles que l'intégration d'un annuaire d'entreprise, la synchronisation avec un ERP, la signature électronique des bons d'affectation ou la mise en production sur une infrastructure cloud complète. Ces extensions constituent néanmoins des perspectives réalistes.

\section{Méthodologie adoptée}
Le projet a été conduit selon une démarche incrémentale inspirée des bonnes pratiques de l'ingénierie logicielle décrites par Pressman \cite{pressman2014} :
\begin{enumerate}
  \item analyse du contexte et cadrage des besoins ;
  \item modélisation fonctionnelle et technique ;
  \item implémentation progressive des couches applicatives ;
  \item validation manuelle et technique des scénarios principaux ;
  \item amélioration continue de l'interface, de la sécurité et de la structure de code.
\end{enumerate}

Cette approche a permis de consolider la cohérence entre les besoins exprimés, les choix d'architecture et les résultats visibles dans l'application finale.

\section{Méthode Agile / Scrum}
Afin de maintenir un rythme de progression régulier et de prioriser les fonctionnalités à forte valeur métier, le projet a été structuré autour d'une logique Agile de type Scrum \cite{schwaber2020}. Le backlog fonctionnel a été organisé en modules, puis découpé en incréments permettant de livrer successivement :
\begin{itemize}
  \item l'ossature technique du projet ;
  \item l'authentification et le contrôle d'accès ;
  \item le cœur métier des équipements et affectations ;
  \item les services complémentaires : maintenance, notifications, rapports ;
  \item l'amélioration de l'interface et de la documentation.
\end{itemize}

Le tableau \ref{tab:scrum-cadence} résume cette logique de progression.

\begin{table}[H]
\centering
\caption{Cadence de réalisation inspirée de Scrum}
\label{tab:scrum-cadence}
\begin{tabularx}{\textwidth}{>{\bfseries}p{2cm}p{3.8cm}X}
\toprule
Sprint & Objectif principal & Livrables dominants \\
\midrule
Sprint 1 & Cadrage et architecture & choix technologiques, structure du dépôt, configuration initiale \\
Sprint 2 & Authentification et sécurité & JWT, rôles, filtres, gestion de session, écrans d'accès \\
Sprint 3 & Inventaire et affectations & CRUD équipements, cycle d'affectation, tableaux métier \\
Sprint 4 & Maintenance et communication & tickets, board, notifications temps réel, historique \\
Sprint 5 & Industrialisation & rapports, Docker, documentation, polish UI/UX \\
\bottomrule
\end{tabularx}
\end{table}

\section{Synthèse du chapitre}
Ce premier chapitre a permis de poser le cadre du projet en mettant en évidence l'enjeu organisationnel de la gestion du parc informatique. Il a montré que la solution attendue devait conjuguer centralisation, sécurité, traçabilité, ergonomie et évolutivité. Le chapitre suivant approfondit cette base en détaillant l'analyse des besoins fonctionnels et non fonctionnels qui ont servi de référence à la conception de la plateforme.









